\documentclass[10pt]{article}

\usepackage[a4paper,left=25mm,right=25mm,top=22mm,bottom=22mm,headheight=14pt,headsep=7mm,footskip=11mm]{geometry}
\usepackage{fontspec}
\usepackage{mathtools}
\usepackage{unicode-math}
\setmainfont{Latin Modern Roman}
\setsansfont{Latin Modern Sans}
\setmathfont{Latin Modern Math}

\usepackage{xcolor}
\definecolor{proofred}{HTML}{7A1825}
\definecolor{proofgray}{HTML}{F4F3F3}
\definecolor{muted}{HTML}{666666}
\usepackage{microtype}
\usepackage{enumitem}
\usepackage{titlesec}
\usepackage{fancyhdr}
\usepackage[most]{tcolorbox}
\usepackage{hyperref}

\hypersetup{
  colorlinks=true,
  linkcolor=proofred,
  urlcolor=proofred,
  pdftitle={Semantic Specifications and Sound Constructions for Summary Reuse},
  pdfsubject={Chapter 5 compact proof note},
  pdfauthor={Mamy}
}

\setlength{\parindent}{0pt}
\setlength{\parskip}{0.34em}
\emergencystretch=1.2em
\setlist[itemize]{leftmargin=1.45em,itemsep=0.2em,topsep=0.2em}

\titleformat{\section}
  {\sffamily\bfseries\color{proofred}\fontsize{14}{17}\selectfont}
  {\thesection}{0.55em}{}
\titlespacing*{\section}{0pt}{1.0em}{0.4em}
\titleformat{\subsection}
  {\sffamily\bfseries\color{proofred}\fontsize{11.5}{14}\selectfont}
  {\thesubsection}{0.5em}{}
\titlespacing*{\subsection}{0pt}{0.72em}{0.22em}

\pagestyle{fancy}
\fancyhf{}
\fancyfoot[C]{\small\thepage}
\renewcommand{\headrulewidth}{0pt}
\renewcommand{\footrulewidth}{0pt}

\newtcolorbox{statementbox}[1][]{
  enhanced,
  breakable,
  colback=proofgray,
  colframe=proofred,
  boxrule=0pt,
  leftrule=2.5pt,
  arc=1.2pt,
  left=8pt,right=8pt,top=6pt,bottom=6pt,
  before skip=0.4em,after skip=0.45em,
  #1
}

\newcommand{\boxheading}[1]{\textbf{#1}\quad}
\newcommand{\Field}{\mathsf{Field}}
\newcommand{\State}{\mathsf{State}}
\newcommand{\Semantics}[1]{\mathcal S\mathopen{⟦}#1\mathclose{⟧}}
\newcommand{\CollectingSemantics}[1]{\widehat{\mathcal S}\mathopen{⟦}#1\mathclose{⟧}}
\newcommand{\AbstractSemantics}[1]{\mathcal S^\sharp\mathopen{⟦}#1\mathclose{⟧}}
\newcommand{\ProjClose}{\mathsf{closure}}
\newcommand{\filter}{\mathsf{filter}}
\newcommand{\relate}{\mathsf{relate}}
\newcommand{\reuse}{\mathsf{coarse\ reuse}}
\newcommand{\frameReuse}{\mathsf{frame\text{-}preserving\ reuse}}
\newcommand{\Lookup}{\mathsf{lookup}}
\newcommand{\SemanticExpand}{\mathsf{semantic\ expansion}}
\newcommand{\SoundFilter}{\mathsf{sound\ filter}}
\newcommand{\ExactFilter}{\mathsf{exact\ filter}}
\newcommand{\SoundReuse}{\mathsf{sound\ reuse}}
\newcommand{\agree}[1]{\mathrel{\equiv_{#1}}}
\newcommand{\qedbox}{\hfill\raisebox{0.15ex}{\fboxsep=0pt\fbox{\rule{0pt}{0.55em}\rule{0.55em}{0pt}}}}

\begin{document}

{\sffamily\bfseries\color{proofred}\fontsize{10.5}{13}\selectfont CHAPTER 5 COMPACT PROPOSITION\par}
\vspace{4mm}
{\sffamily\bfseries\fontsize{23}{27}\selectfont
Semantic Specifications and\\
Sound Constructions for Summary Reuse\par}
\vspace{4mm}
{\color{muted}\large Replay sets, exact filtering, summary production, and sound reuse\par}
\vspace{4mm}
{\color{proofred}\rule{\textwidth}{1pt}\par}

\section{Setting and terminology}

Fix a statement or function body \(s\). Let \(\Field\) denote the set of fields and
\(\State\) the set of concrete states. Let \(\mathbb D^\sharp\) be an abstract domain, with
\[
  \gamma:\mathbb D^\sharp\to\mathcal P(\State)
\]
as its concretization. The concrete semantics of \(s\) is a relation between input and output
states; its collecting semantics maps a set of input states to the possible output states:
\[
  \Semantics{s}\subseteq\State\times\State,
  \qquad
  \CollectingSemantics{s}:\mathcal P(\State)\to\mathcal P(\State),
\]
where
\[
  \CollectingSemantics{s}(X)
  =\{O\in\State\mid \exists I\in X,\ (I,O)\in\Semantics{s}\}.
\]

We call the analysis that produces a summary the \emph{producer}. It analyzes \(s\) from an
abstract input \(I^\sharp\), obtains \(O^\sharp\), and stores the cache entry
\[
  c=(A,W,I_c^\sharp,O_c^\sharp),
\]
where \(A\) is the \emph{replay set}, \(W\) is a sound write set, and
\(I_c^\sharp,O_c^\sharp\) are the stored input and output. We call a later analysis of the
same \(s\), starting from the current abstract input \(J^\sharp\), the \emph{consumer}.
\emph{Lookup}, also called \emph{matching}, checks whether the consumer input and the stored
input agree on the replay set \(A\). When lookup succeeds, \emph{replay}, also called
\emph{reuse}, uses the stored summary result instead of reanalyzing \(s\) for the consumer.
\emph{Coarse reuse} returns only \(O_c^\sharp\). Consequently, coarse reuse does not retain
information from the consumer input \(J^\sharp\), including information about fields outside
\(A\cup W\). By contrast, \emph{frame-preserving reuse} also retains the consumer's
information about fields outside the write set \(W\).

Reuse is sound only when it preserves the input fields in the replay set \(A\). A sound
\emph{read set} \(R\) contains every input field whose value \(s\) may read or otherwise
depend on; a sound write set \(W\) contains every field that \(s\) may modify.
The operation \(\relate^\sharp\) adds any field that must be retained together with the selected
fields to preserve a relation represented in the abstract state. The operation \(\filter^\sharp\) retains information about the selected fields and
discards information about all other fields. Their types and the semantic properties they must
satisfy are specified below.

For \(A\subseteq\Field\), write \(I\agree{A}J\) when the concrete states \(I\) and \(J\)
agree on every field in \(A\). The agreement closure maps a set of fields and a set of concrete
states to a set of concrete states:
\[
  \ProjClose:\mathcal P(\Field)\times\mathcal P(\State)
  \to\mathcal P(\State)
\]
and is defined by
\[
  \ProjClose(A,X)=\{J\in\State\mid \exists I\in X,\ I\agree{A}J\}.
\]

\section{Semantic specifications}

\subsection{Replay set specification}

Semantic expansion specifies when a replay set \(A\) is sufficient to preserve the possible
outputs on fields \(B\). It is the predicate
\[
  \SemanticExpand:
  \mathcal P(\Field)\times\mathcal P(\Field)\times\mathcal P(\State)
  \to\mathsf{Prop},
\]
defined by
\[
  \SemanticExpand(A;B,P)
  \Longleftrightarrow
  \CollectingSemantics{s}\!\left(\ProjClose(A,P)\right)
  \subseteq
  \ProjClose\!\left(B,\CollectingSemantics{s}(P)\right).
  \tag{SE}
\]
Thus, \(A\) is sufficient for replay on the output fields \(B\): varying the values outside
\(A\), while keeping \(A\) fixed, cannot create a new possible output on \(B\).

\subsection{Filter specifications}

Sound and exact filtering both take three arguments: \(B\), the fields to retain;
\(S^\sharp\), the state before filtering; and \(T^\sharp\), the filtered state:
\[
  \SoundFilter,\ExactFilter:
  \mathcal P(\Field)\times\mathbb D^\sharp\times\mathbb D^\sharp
  \to\mathsf{Prop},
\]
defined by
\begin{align*}
  \SoundFilter(B,S^\sharp,T^\sharp)
  &\Longleftrightarrow
  \ProjClose\!\left(B,\gamma(S^\sharp)\right)\subseteq\gamma(T^\sharp),\\
  \ExactFilter(B,S^\sharp,T^\sharp)
  &\Longleftrightarrow
  \gamma(T^\sharp)=\ProjClose\!\left(B,\gamma(S^\sharp)\right).
\end{align*}
\subsection{Reuse specification}

Sound reuse takes two arguments: \(J^\sharp\), the consumer input, and \(Q^\sharp\), the
reuse result:
\[
  \SoundReuse:\mathbb D^\sharp\times\mathbb D^\sharp\to\mathsf{Prop},
\]
defined by
\[
  \SoundReuse(J^\sharp,Q^\sharp)
  \Longleftrightarrow
  \CollectingSemantics{s}\!\left(\gamma(J^\sharp)\right)
  \subseteq\gamma(Q^\sharp).
  \tag{SR}
\]

\section{Replay set construction}

Let \(R\) and \(W\) be sound read and write sets for \(s\), and let \(I^\sharp\) be the
producer input. The operation \(\relate^\sharp\) extends a candidate replay set with fields
needed to preserve relations in \(I^\sharp\):
\[
  \relate^\sharp:
  \mathcal P(\Field)\times\mathbb D^\sharp\to\mathcal P(\Field).
\]

A replay set \(A\) must satisfy the following \emph{Replay Closure} conditions:
\[
  R\subseteq A,
  \qquad
  \relate^\sharp(A,I^\sharp)\subseteq A.
  \tag{Replay Closure}
\]
The first condition retains every field that \(s\) may read. The second retains every field
needed to preserve relations including the fields already selected. Together, they ensure that
retaining the fields in \(A\) is enough to reproduce the possible outputs of \(s\) on \(A\cup W\),
and let us obtain an exact projection of the relational producer state.

To construct a smaller replay set when only finitely many fields can be added, start from the
read set and repeatedly add relational dependencies:
\[
  A_0=R,
  \qquad
  A_{n+1}=A_n\cup\relate^\sharp(A_n,I^\sharp),
\]
until the sequence stabilizes. The result satisfies Replay Closure. A replay set always exists,
because \(A=\Field\) satisfies the two closure conditions. The proofs below require only
Replay Closure, not termination or leastness.

We assume that the read and write sets satisfy the following semantic property:
\[
  R\subseteq A
  \quad\Longrightarrow\quad
  \SemanticExpand(A;A\cup W,P)
  \quad\text{for every }P\subseteq\State.
  \tag{Read/Write Expansion}
\]

\begin{statementbox}
\boxheading{Proposition 1 --- Replay set correctness}
Under Replay Closure, fixing the fields in \(A\) is enough to determine the possible outputs on
\(A\cup W\):
\[
  \SemanticExpand(A;A\cup W,\gamma(I^\sharp)).
\]
\end{statementbox}

\textbf{Proof.}
Replay Closure gives \(R\subseteq A\); Read/Write Expansion gives the conclusion. \qedbox

\section{Filter construction}

The abstract domain provides a filtering operation that retains information about a selected set
of fields:
\[
  \filter^\sharp:
  \mathcal P(\Field)\times\mathbb D^\sharp\to\mathbb D^\sharp.
\]
It must satisfy the following two properties:
\begin{align*}
  &\SoundFilter\!\left(B,S^\sharp,\filter^\sharp(B,S^\sharp)\right),
  &&\tag{Filter Sound}\\
  &\relate^\sharp(B,S^\sharp)\subseteq B
  \Longrightarrow
  \ExactFilter\!\left(B,S^\sharp,\filter^\sharp(B,S^\sharp)\right).
  &&\tag{Closed Exact}
\end{align*}

Filter Sound requires the filtered state to safely cover the retained information. Closed Exact
requires it to represent that information exactly when the selected fields already contain all
their relational dependencies.

The following equation expresses when an abstract state separates exactly across a field set
\(A\) and its complement:
\[
  \gamma(S^\sharp)
  =
  \ProjClose\!\left(A,\gamma(S^\sharp)\right)
  \cap
  \ProjClose\!\left(\Field\setminus A,\gamma(S^\sharp)\right).
  \tag{Separation}
\]
When this equation holds, the information represented by \(S^\sharp\) contains no relation
crossing between \(A\) and its complement. It gives the semantic intuition behind relation
closure.

\begin{statementbox}
\boxheading{Proposition 2 --- Exact filtering of the producer input}
Under Replay Closure, filtering the producer input to \(A\) is exact:
\[
  \ExactFilter\!\left(A,I^\sharp,\filter^\sharp(A,I^\sharp)\right).
\]
\end{statementbox}

\textbf{Proof.}
Replay Closure gives \(\relate^\sharp(A,I^\sharp)\subseteq A\); Closed Exact gives the
conclusion. \qedbox

Production analyzes \(I^\sharp\) before filtering it. Therefore,
\(\filter^\sharp(A,I^\sharp)\) must represent \(\ProjClose(A,\gamma(I^\sharp))\) exactly,
and this agreement closure must be representable in \(\mathbb D^\sharp\). By contrast, sound
filtering suffices for the stored output and lookup.

\section{Summary production}

The abstract semantics maps each abstract state to an abstract state and soundly overapproximates
its collecting semantics:
\[
  \AbstractSemantics{s}:\mathbb D^\sharp\to\mathbb D^\sharp
\]
and satisfies analyzer soundness:
\[
  \CollectingSemantics{s}\!\left(\gamma(S^\sharp)\right)
  \subseteq\gamma\!\left(\AbstractSemantics{s}(S^\sharp)\right).
  \tag{Analyze Sound}
\]
For the producer input \(I^\sharp\), define
\[
  O^\sharp=\AbstractSemantics{s}(I^\sharp),
  \qquad
  I_c^\sharp=\filter^\sharp(A,I^\sharp),
  \qquad
  O_c^\sharp=\filter^\sharp(A\cup W,O^\sharp),
\]
and store the summary
\[
  c=(A,W,I_c^\sharp,O_c^\sharp).
\]

\begin{statementbox}
\boxheading{Proposition 3 --- Summary soundness}
The stored output safely covers the possible outputs from the stored input:
\[
  \CollectingSemantics{s}\!\left(\gamma(I_c^\sharp)\right)
  \subseteq\gamma(O_c^\sharp).
  \tag{Summary Sound}
\]
\end{statementbox}

\textbf{Proof.}
Propositions 1 and 2, analyzer soundness, monotonicity of \(\ProjClose(A\cup W,\cdot)\), and
Filter Sound give
\begin{align*}
  \CollectingSemantics{s}\!\left(\gamma(I_c^\sharp)\right)
  &=\CollectingSemantics{s}\!\left(\ProjClose(A,\gamma(I^\sharp))\right)\\
  &\subseteq\ProjClose\!\left(A\cup W,
      \CollectingSemantics{s}(\gamma(I^\sharp))\right)\\
  &\subseteq\ProjClose\!\left(A\cup W,\gamma(O^\sharp)\right)\\
  &\subseteq\gamma\!\left(\filter^\sharp(A\cup W,O^\sharp)\right)
   =\gamma(O_c^\sharp).
\end{align*}
\qedbox

\section{Sound coarse reuse}

Lookup for a summary \(c\) succeeds exactly when the consumer input, filtered to \(A\), equals
the stored input:
\[
  \Lookup:\mathbb D^\sharp\to\mathsf{Prop},
  \qquad
  \Lookup(J^\sharp)
  \Longleftrightarrow
  \filter^\sharp(A,J^\sharp)=I_c^\sharp.
\]
When lookup succeeds, coarse reuse returns the stored output \(O_c^\sharp\):
\[
  \reuse^\sharp:\mathbb D^\sharp\rightharpoonup\mathbb D^\sharp,
  \qquad
  \reuse^\sharp(J^\sharp)=O_c^\sharp
  \quad\text{when }\Lookup(J^\sharp).
  \tag{Coarse Reuse}
\]

Filter Sound ensures that every successful lookup satisfies the following inclusion:
\[
  \Lookup(J^\sharp)
  \Longrightarrow
  \gamma(J^\sharp)\subseteq\gamma(I_c^\sharp).
  \tag{Match Sound}
\]
Indeed, \(\gamma(J^\sharp)\subseteq
\ProjClose(A,\gamma(J^\sharp))\subseteq
\gamma(\filter^\sharp(A,J^\sharp))=\gamma(I_c^\sharp)\).

\begin{statementbox}
\boxheading{Proposition 4 --- Sound summary reuse}
A successful coarse reuse safely covers consumer outputs:
\[
  \Lookup(J^\sharp)
  \quad\Longrightarrow\quad
  \SoundReuse\!\left(J^\sharp,\reuse^\sharp(J^\sharp)\right).
\]
\end{statementbox}

\textbf{Proof.}
Match Sound, monotonicity of \(\CollectingSemantics{s}\), Summary Sound, and Coarse Reuse give
\[
  \CollectingSemantics{s}\!\left(\gamma(J^\sharp)\right)
  \subseteq
  \CollectingSemantics{s}\!\left(\gamma(I_c^\sharp)\right)
  \subseteq\gamma(O_c^\sharp)
  =\gamma\!\left(\reuse^\sharp(J^\sharp)\right).
\]
This is exactly \(\SoundReuse(J^\sharp,\reuse^\sharp(J^\sharp))\). \qedbox

\section{Sound frame-preserving reuse}

To retain consumer information about fields that \(s\) does not write, assume the following
Write Frame property:
\[
  (J,O)\in\Semantics{s}
  \quad\Longrightarrow\quad
  J\agree{\Field\setminus W}O.
  \tag{Write Frame}
\]
The abstract domain also provides a meet operation:
\[
  \sqcap^\sharp:\mathbb D^\sharp\times\mathbb D^\sharp\to\mathbb D^\sharp
\]
It must satisfy two separate properties:
\begin{align*}
  \gamma(S_1^\sharp)\cap\gamma(S_2^\sharp)
  &\subseteq\gamma(S_1^\sharp\sqcap^\sharp S_2^\sharp),
  &&\tag{Meet Sound}\\
  \gamma(S_1^\sharp\sqcap^\sharp S_2^\sharp)
  &\subseteq\gamma(S_1^\sharp)\cap\gamma(S_2^\sharp).
  &&\tag{Meet Reductive}
\end{align*}

Together, these properties make the meet represent the intersection of the two concretizations
exactly. When lookup succeeds, frame-preserving reuse intersects the stored output with the
consumer information on fields outside \(W\):
\[
  \frameReuse^\sharp(J^\sharp)
  =O_c^\sharp\sqcap^\sharp
    \filter^\sharp(\Field\setminus W,J^\sharp).
  \tag{Frame-Preserving Reuse}
\]

\begin{statementbox}
\boxheading{Proposition 5 --- Sound and precise frame-preserving reuse}
When lookup succeeds, frame-preserving reuse is sound and at least as precise as coarse reuse:
\begin{align*}
  &\text{Sound:}
  &&\SoundReuse\!\left(J^\sharp,\frameReuse^\sharp(J^\sharp)\right),\\
  &\text{Precise:}
  &&\gamma\!\left(\frameReuse^\sharp(J^\sharp)\right)
    \subseteq\gamma\!\left(\reuse^\sharp(J^\sharp)\right).
\end{align*}
\end{statementbox}

\textbf{Proof of Sound.}
Take \(J\in\gamma(J^\sharp)\) and \((J,O)\in\Semantics{s}\). Proposition~4 gives
\(O\in\gamma(O_c^\sharp)\). Write Frame gives
\(J\agree{\Field\setminus W}O\), hence
\[
  O\in\ProjClose\!\left(\Field\setminus W,\gamma(J^\sharp)\right).
\]
Filter Sound then gives
\[
  O\in\gamma\!\left(\filter^\sharp(\Field\setminus W,J^\sharp)\right).
\]
Thus \(O\) belongs to the intersection of the concretizations of both operands in
Frame-Preserving Reuse. Meet Sound gives
\(O\in\gamma(\frameReuse^\sharp(J^\sharp))\). \qedbox

\textbf{Proof of Precise.}
Meet Reductive and Coarse Reuse give
\[
  \gamma\!\left(\frameReuse^\sharp(J^\sharp)\right)
  \subseteq\gamma(O_c^\sharp)
  =\gamma\!\left(\reuse^\sharp(J^\sharp)\right).
\]
Therefore, frame-preserving reuse admits no more concrete states than coarse reuse. \qedbox

\end{document}
